\section{Architectural Tactics}

The tactics address how to adjust a deployment's position in the design space by targeting its autonomy, agency, or both. They operate at three layers---workflow, tool, and action---and can raise or lower the targeted dimension depending on the deployment goal. Table~\ref{tab:tactics} summarises all six.

\begin{table}[t]
    \centering
    \caption{Six architectural tactics for navigating the autonomy--agency design space. \emph{Axis}: the design dimension primarily targeted. \emph{Layer}: where the tactic operates --- workflow (control flow), tool (available capabilities), or action (reversibility). \emph{Effect}: whether the tactic raises~($\uparrow$), lowers~($\downarrow$), or can do either~($\uparrow\downarrow$) the target dimension.}\label{tab:tactics}
    \footnotesize
    \setlength{\tabcolsep}{5pt}
    \begin{tabular}{llllc}
        \toprule
        \textbf{ID} & \textbf{Tactic} & \textbf{Axis} & \textbf{Layer} & \textbf{Effect} \\
        \midrule
        T1 & Checkpoints            & Autonomy & Workflow & $\downarrow$ \\[2pt]
        T2 & Escalation             & Autonomy & Workflow & $\downarrow$ \\[2pt]
        T3 & Multi-Agent Delegation & Autonomy & Workflow & $\uparrow$   \\
        \midrule
        T4 & Tool Provisioning      & Agency   & Tool     & $\uparrow\downarrow$ \\[2pt]
        T5 & Tool Fencing           & Agency   & Tool     & $\downarrow$ \\[2pt]
        T6 & Write Staging          & Agency   & Action   & $\downarrow$ \\
        \bottomrule
    \end{tabular}
\end{table}

\paragraph{Autonomy tactics.}
\textbf{Checkpoints}~(T1) insert fixed, unconditional pauses at predetermined workflow steps, making review locations auditable by design. \textbf{Escalation}~(T2) is a dynamic, agent-initiated pause that fires when the agent determines it cannot proceed autonomously---preserving high autonomy in the common case while restoring human involvement at the edge of the agent's reliable operating boundary. \textbf{Multi-agent delegation}~(T3) raises effective autonomy by replacing a human sub-task handler with a specialised sub-agent; the delegation boundary introduces a gap in the audit trail that requires explicit design attention.

\paragraph{Agency tactics.}
\textbf{Tool provisioning}~(T4) adjusts which tools the agent can access at runtime. \textbf{Tool fencing}~(T5) constrains a tool's operational envelope without removing it, implemented in the tool's own validation layer independently of model behaviour. \textbf{Write staging}~(T6) interposes a draft layer before an authoritative commit, demoting an irreversible L5 write to a reversible L4 action---converting a permanent write into one that can be reviewed and corrected before it takes effect.

